\documentclass[11pt]{article}
\usepackage[hyperref]{acl}
\usepackage{times}
\usepackage{latexsym}
\usepackage{microtype}

\title{UnMask: Supplementary Material\\
\large Conversation Transcripts (Appendix A)}

\author{Sanika Vilas Najan \quad Vaishak Girish Kumar \\
  University at Buffalo \\
  CSE 635: NLP and Text Mining --- Spring 2026}

\begin{document}
\maketitle

\section*{Appendix A: Conversation Transcripts}
\label{app:transcripts}

The P3 specification requires 5 conversation transcripts demonstrating no
answer leak before the 3rd turn.
All transcripts were generated live from UnMask with cold-start mastery = 0.20
(PCR \texttt{context\_only} mode).
Gold answers are shown for grader reference only---they were present in
\texttt{InternalAnalysis} (hidden) and absent from \texttt{VisibleResponse}
(shown to student).
Leak status verified by both layers of the answer-leak detector.

\subsection*{Transcript 1 --- Brachial Plexus Origins}

\noindent\textbf{Turn 1} $|$ PCR: \texttt{context\_only} $|$ Leak: None $|$ Ends `?': Yes\\
\textit{Student:} What spinal cord levels form the brachial plexus?\\
\textit{UnMask:} Great question! Let's explore the spinal levels involved.
Which specific cervical and thoracic spinal nerves contribute to the formation
of the brachial plexus, and how do they relate to the vertebrae they emerge
from?\\
{\small\textit{[Gold held internally: ``C5--T1 anterior rami'']}}

\noindent\textbf{Turn 2} $|$ PCR: \texttt{context\_only} $|$ Leak: None $|$ Ends `?': Yes\\
\textit{Student:} Is it only cervical nerves?\\
\textit{UnMask:} You're on the right track! Have you considered whether any
thoracic spinal nerves might also contribute to the brachial plexus, in
addition to the cervical nerves?\\
{\small\textit{[Gold held internally: ``Lower 4 cervical (C5--C8) and first thoracic (T1)'']}}

\noindent\textbf{Turn 3} $|$ PCR: \texttt{prerequisite\_first} $|$ Leak: None $|$ Ends `?': Yes\\
\textit{Student:} So five roots total from C5 to T1?\\
\textit{UnMask:} You're doing well in identifying the spinal roots involved!
How do these spinal nerve roots organize themselves into the trunks and
divisions of the brachial plexus?\\
{\small\textit{[Gold confirmed: ``Yes---C5, C6, C7, C8, T1'']}}

\subsection*{Transcript 2 --- Axillary Nerve \& Shoulder Dislocation}

\noindent\textbf{Turn 1} $|$ PCR: \texttt{context\_only} $|$ Leak: None $|$ Ends `?': Yes\\
\textit{Student:} Which nerve is damaged when a patient can't abduct past 15 degrees after shoulder dislocation?\\
\textit{UnMask:} What is the primary muscle responsible for shoulder abduction beyond 15 degrees, and which nerve innervates this muscle?\\
{\small\textit{[Gold held: ``Axillary nerve'']}}

\noindent\textbf{Turn 2} $|$ PCR: \texttt{context\_only} $|$ Leak: None $|$ Ends `?': Yes\\
\textit{Student:} So is it the axillary nerve?\\
\textit{UnMask:} You're on the right track! What is the primary muscle responsible for shoulder abduction beyond 15 degrees, and which nerve innervates this muscle?\\
{\small\textit{[Gold held: ``Axillary nerve innervates deltoid, crucial for abduction beyond 15°'']}}

\noindent\textbf{Turn 3} $|$ PCR: \texttt{prerequisite\_first} $|$ Leak: None $|$ Ends `?': Yes\\
\textit{Student:} What else does the axillary nerve innervate besides deltoid?\\
\textit{UnMask:} Besides the deltoid, which other muscle and what specific area of skin does the axillary nerve innervate?\\
{\small\textit{[Gold confirmed: ``Teres minor + lateral shoulder skin (regimental badge area)'']}}

\subsection*{Transcript 3 --- Radial Nerve \& Wrist Drop}

\noindent\textbf{Turn 1} $|$ PCR: \texttt{context\_only} $|$ Leak: None $|$ Ends `?': Yes\\
\textit{Student:} A patient has wrist drop after a mid-shaft humerus fracture. What nerve is likely injured?\\
\textit{UnMask:} What is the relationship between the location of a mid-shaft humerus fracture and the nerve that runs in that area, particularly in relation to wrist extensor function?\\
{\small\textit{[Gold held: ``Radial nerve'']}}

\noindent\textbf{Turn 2} $|$ PCR: \texttt{context\_only} $|$ Leak: None $|$ Ends `?': Yes\\
\textit{Student:} Does the radial nerve innervate all wrist extensors?\\
\textit{UnMask:} What is the role of the radial nerve in the forearm, and how does it relate to the muscles responsible for wrist extension?\\
{\small\textit{[Gold held: ``Radial nerve innervates all wrist extensor muscles'']}}

\noindent\textbf{Turn 3} $|$ PCR: \texttt{prerequisite\_first} $|$ Leak: None $|$ Ends `?': Yes\\
\textit{Student:} What about sensation---where would the patient feel numbness?\\
\textit{UnMask:} Which specific area on the dorsal side of the hand and forearm would you expect to experience numbness if the radial nerve is injured?\\
{\small\textit{[Gold confirmed: ``Posterior forearm and dorsal hand; first dorsal web space'']}}

\subsection*{Transcript 4 --- Rotator Cuff Anatomy}

\noindent\textbf{Turn 1} $|$ PCR: \texttt{context\_only} $|$ Leak: None $|$ Ends `?': Yes\\
\textit{Student:} Name the four rotator cuff muscles.\\
\textit{UnMask:} What acronym can help you remember the names of the four muscles that form the rotator cuff, and what does each letter stand for?\\
{\small\textit{[Gold held: ``Supraspinatus, infraspinatus, teres minor, subscapularis (SITS)'']}}

\noindent\textbf{Turn 2} $|$ PCR: \texttt{context\_only} $|$ Leak: None $|$ Ends `?': Yes\\
\textit{Student:} I know supraspinatus---what about the others on the posterior scapula?\\
\textit{UnMask:} What are the names of the rotator cuff muscles that originate from the infraspinous fossa and the teres minor region on the posterior scapula?\\
{\small\textit{[Gold held: ``Infraspinatus and teres minor'']}}

\noindent\textbf{Turn 3} $|$ PCR: \texttt{prerequisite\_first} $|$ Leak: None $|$ Ends `?': Yes\\
\textit{Student:} What movement does the infraspinatus produce?\\
\textit{UnMask:} Considering the infraspinatus is located behind the glenohumeral joint, how might its position influence the type of rotation it facilitates at the shoulder?\\
{\small\textit{[Gold confirmed: ``External rotation of the shoulder'']}}

\subsection*{Transcript 5 --- Median Nerve \& Carpal Tunnel Syndrome}

\noindent\textbf{Turn 1} $|$ PCR: \texttt{context\_only} $|$ Leak: None $|$ Ends `?': Yes\\
\textit{Student:} What nerve is compressed in carpal tunnel syndrome?\\
\textit{UnMask:} Which nerve passes through the carpal tunnel along with the flexor tendons and is responsible for sensation in the thumb, index, middle, and part of the ring finger?\\
{\small\textit{[Gold held: ``Median nerve'']}}

\noindent\textbf{Turn 2} $|$ PCR: \texttt{context\_only} $|$ Leak: None $|$ Ends `?': Yes\\
\textit{Student:} What symptoms would the patient have?\\
\textit{UnMask:} When the median nerve is compressed at the wrist, what changes might you expect in both sensation and muscle function in the hand?\\
{\small\textit{[Gold held: ``Numbness/tingling digits 1--3.5; thenar weakness'']}}

\noindent\textbf{Turn 3} $|$ PCR: \texttt{prerequisite\_first} $|$ Leak: None $|$ Ends `?': Yes\\
\textit{Student:} Would the patient have weakness gripping or just sensory loss?\\
\textit{UnMask:} What role do the thenar muscles, innervated by the median nerve, play in gripping, and how might their dysfunction affect hand strength?\\
{\small\textit{[Gold confirmed: ``Both---thenar weakness impairs pinch/grip; sensory loss in lateral digits'']}}

\bigskip
\noindent\textbf{Transcript summary:} 15 turns, 5 transcripts, 0 confirmed leaks,
0 soft flags, 15/15 ends with ``?''.
Gold answer present only in \texttt{InternalAnalysis} (stripped before
rendering); absent from retrieved Qdrant chunks due to PCR \texttt{must\_not}
filter.

\end{document}
